% Generated draft: requires mathematical and source review.
\documentclass[11pt]{article}
\usepackage[margin=25mm]{geometry}
\usepackage{fontspec}
\setmainfont{Latin Modern Roman}
\usepackage{amsmath,amssymb,mathtools,unicode-math}
\setmathfont{Latin Modern Math}
\usepackage{enumitem,needspace,xurl,hyperref}
\setlength{\parindent}{0pt}
\setlength{\parskip}{5pt}
\setlength{\emergencystretch}{3em}
\newcommand{\N}{\mathbb{N}}
\newcommand{\Z}{\mathbb{Z}}
\newcommand{\Q}{\mathbb{Q}}
\newcommand{\R}{\mathbb{R}}
\newcommand{\C}{\mathbb{C}}
\newcommand{\F}{\mathbb{F}}
\newcommand{\A}{\mathbb{A}}
\newcommand{\PP}{\mathbb{P}}
\newcommand{\E}{\mathbb{E}}
\newcommand{\eps}{\varepsilon}
\newcommand{\dd}{\,\mathrm{d}}
\begin{document}
\section*{UnsolvedMath ID: 20001006}
\textbf{Source problem code:} catalogue entry
\subsection{Definitions and mathematical statement}
\textbf{Item 41 (Lionel Levine).} Study the extra automorphism-group structure on the critical group of a graph carrying a harmonic action. The source is a research suggestion, not a fixed conjecture or requested classification. The registry says open/\texttt{OPEN-TRIAGE}, while the aid labels the item context-only/not-a-problem.

For a finite connected loopless multigraph $X$, let $\operatorname{Div}^0(X)$ be degree-zero integral vertex divisors and
\[
 \operatorname{Jac}(X)=\operatorname{Div}^0(X)/\operatorname{im}\Delta
\]
its critical group. Every graph automorphism acts on vertex divisors and commutes with the Laplacian, inducing $\operatorname{Aut}(X)\to\operatorname{Aut}(\operatorname{Jac}(X))$. A specified group action therefore gives a finite group module. A usual harmonic-action convention, without edge inversions, requires vertex stabilizers to act freely on incident half-edges; the original does not supply its convention.

\textbf{Reported kernel-displacement theorem.} Suppose $X$ is bridgeless with $n$ vertices. Let
\[
 K=\ker(\operatorname{Aut}(X)\to\operatorname{Aut}(\operatorname{Jac}(X))),\qquad
 P=\ker(\operatorname{Aut}(X)\to\operatorname{Sym}(V(X))).
\]
Vertex-fixing edge permutations lie in $K$. The aid proves that
\[
 d(\sigma)=[\sigma(v)-v]\quad(\sigma\in K)
\]
is independent of $v$, is a homomorphism, and induces the canonical injection
\[
 K/P\hookrightarrow\operatorname{Jac}(X)[n].
\]
Here $[n]$ denotes elements annihilated by $n$. For simple bridgeless graphs $P$ is trivial, so the automorphisms invisible on the critical group form an abelian group of exponent dividing $n$. The bridge-free hypothesis is material to identifying the zero-displacement kernel with $P$.

The aid combines this with a 2025 semiregular faithfulness theorem of Estelyi--Karabas--Mednykh--Nedela to conclude that the full automorphism group of every simple connected $3$-edge-connected graph acts faithfully on its critical group. Here $3$-edge-connected means no cut of fewer than three edges disconnects the graph. The exact external theorem hypotheses and citation are absent.

\textbf{Reported multigraph contrast.} For any finite group $G$, take the banana multigraph with two vertices and $|G|$ parallel edges labeled by $G$. Let $G$ fix the vertices and act regularly on edge labels. This graph action is faithful and harmonic in the declared convention, while the induced action on
\[
 \operatorname{Jac}(X)\simeq\Z/|G|\Z
\]
is trivial because all vertex divisors are fixed. Thus a faithful harmonic graph action need not be detectable on the critical group when parallel-edge permutations are allowed. This contrasts with the simple $3$-edge-connected conclusion and does not classify all possible modules or settle a formal conjecture absent from the original suggestion.
\subsection{Short English statement}
The source suggests studying a graph's critical group together with a harmonic automorphism action. The aid gives a canonical bound on the invisible automorphisms of bridgeless graphs and reports faithfulness for simple $3$-edge-connected graphs. Parallel-edge banana examples show that arbitrary faithful harmonic graph actions can still act trivially on the critical group.
\subsection{Sources}
\begin{itemize}
\item Local UnsolvedMath registry, permanent ID $20001006$, \texttt{catalogue entry}; canonical \texttt{aim-combinatorics-notes.json}, \texttt{notes[130]}.
\item AIM workshop ``Generalizations of chip-firing and the critical group,'' PDF item 41, Lionel Levine: \url{https://aimath.org/pastworkshops/chipfiringproblems.pdf}.
\item Estelyi--Karabas--Mednykh--Nedela, 2025 semiregular faithfulness theorem, as named in the supplied summary; no full title, identifier or theorem number supplied.
\item Research attempt 1, \texttt{context\_only}, run status \texttt{not\_a\_problem}, internal rigor \texttt{proved}, low novelty confidence. Explicitly machine-generated, not peer reviewed, and to be independently verified before citation. Original literature field empty.
\end{itemize}
\end{document}
